% !TEX TS-program = pdflatex
% !TEX encoding = UTF-8 Unicode
% !TEX spellcheck = es_CL
% !TEX root = Plantilla_PMT_MT.tex

%\usepackage{amsmath}
\usepackage{mathtools}
\usepackage{amsfonts}
\usepackage{amssymb}
\usepackage{fancyhdr}

\usepackage{xcolor}

\usepackage{graphicx}

\usepackage[hypcap=true]{caption}
%\captionsetup{labelfont=bf} % solo etiqueta en negrita
\captionsetup{font={bf,small}}% todo el titulo en negrita
\usepackage[hypcap=true]{subcaption} % reemplaza a subfig y subfloats

\usepackage{longtable}

\usepackage{booktabs}

%\usepackage[numbers]{natbib}

% para citar y mostrar el texto completo.
% El paquete IEEEtrantools y una entrada
% adicional en la base de datos de la
% bibliografia soluciona los problemas que
% se producen cuando las ref tienen url
\usepackage[retainorgcmds]{IEEEtrantools}
\usepackage{bibentry} % \bibentry
\nobibliography*

% redefinicion del comando bibentry para que los enlaces 
% de las citas apunten a la bibliografia y no a donde se
% usa el bibentry
\makeatletter
\renewcommand\bibentry[1]{\nocite{#1}{\frenchspacing
     \@nameuse{BR@r@#1\@extra@b@citeb}}}
%% original de bibentry.sty
%\newcommand\bibentry[1]{\nocite{#1}{\frenchspacing
%     \hyper@natanchorstart{#1\@extra@b@citeb}%
%     \@nameuse{BR@r@#1\@extra@b@citeb}\hyper@natanchorend}}
\makeatother

%  para cambiar interlineado
% en diferentes partes del texto
% \setstretch{1.3}
\usepackage{setspace}
% Salto de parrafo, genera mas espacio en el indice tambien
%\setlength{\parskip}{0.4cm}

% Agregar directorio donde buscar figuras
\graphicspath{{./figuras/}}

\definecolor{tableau10_1}{RGB}{31,119,180}
\definecolor{tableau10_2}{RGB}{255,127,14}
\definecolor{tableau10_3}{RGB}{44,160,44}
\definecolor{tableau10_4}{RGB}{214,39,40}
\definecolor{tableau10_5}{RGB}{148,103,189}
\definecolor{tableau10_6}{RGB}{140,86,75}
\definecolor{tableau10_7}{RGB}{227,119,194}
\definecolor{tableau10_8}{RGB}{127,127,127}
\definecolor{tableau10_9}{RGB}{188,189,34}
\definecolor{tableau10_10}{RGB}{23,190,207}
\definecolor{gray75}{gray}{0.75}

\definecolor{MSBlue}{rgb}{.204,.353,.541}
\definecolor{MSLightBlue}{rgb}{.31,.506,.741}

%\usepackage{sectsty}
%\chapterfont{\Large}
%\chapterfont{\color{tableau10_1}}
%\sectionfont{\color{tableau10_1}}
%\allsectionsfont{\color{tableau10_1}} % equivalente a arriba

%%%%%%%%%%%%%%%%%%%%%%%%%%%%%%%%%%%%%%%%%%%%%%%%%%%%%%
%% PERSONALIZACION DE TITULOS
%%%%%%%%%%%%%%%%%%%%%%%%%%%%%%%%%%%%%%%%%%%%%%%%%%%%%%

\newcommand{\hsp}{\hspace{20pt}}
\newcommand{\chaprule}{\titlerule[1.5pt]}

%\usepackage{titlesec}
%% Formato de titulos oficial
%\titleformat{\chapter}[hang] % Customise the \section command
%%{\Huge\scshape\raggedright}%
%%{\LARGE\sffamily\bfseries\raggedright}%
%{\LARGE\bfseries\raggedright}%
%{\chaptertitlename~\thechapter.}{1em} % Can be used to give a prefix to all sections, like 'Section ...'
%{} % Can be used to insert code before the heading
%[\chaprule] % Inserts a horizontal line after the heading 
%\titlespacing*{\chapter}{0pt}{-50pt}{40pt}

%\titleformat{\chapter}[hang]{\color{tableau10_1}\Huge\bfseries}{\thechapter\hsp\textcolor{gray75}{|}\hsp}{0pt}{\color{tableau10_1}\Huge\bfseries}
%\titleformat*{\section}{\color{tableau10_1}}

\fancyhf{}
\fancyhead[l]{\footnotesize \textsc{Universidad de Concepción}\\
\textsc{Facultad de Ingeniería}\\
\textsc{Departamento de Ingeniería Eléctrica}}
\fancyhead[r]{\includegraphics[scale=0.12]{logoudec}}
\fancyfoot[r]{\thepage}
\fancyfoot[c]{}
%\pagestyle{fancy}

%\usepackage{transparent}
%\usepackage{draftwatermark}
%\SetWatermarkAngle{0}
%\SetWatermarkText{\transparent{0.1}{\includegraphics{logoudec}}}

% activar solo en la version final
% para que la compilacion no tome tanto tiempo
%\usepackage[pages=some]{background}
\usepackage{background}

\backgroundsetup{
scale=1,
color=black,
opacity=0.05,
angle=0,
contents={%
  \includegraphics[scale=0.5]{logoudec}
  }%
}

\newcommand{\monthname}{%
  \ifcase\the\month
  \or Enero% 1
  \or Febrero% 2
  \or Marzo% 3
  \or Abril% 4
  \or Mayo% 5
  \or Junio% 6
  \or Julio% 7
  \or Agosto% 8
  \or Septiembre% 9
  \or Octubre% 10
  \or Noviembre% 11
  \or Diciembre% 12
  \fi}

%\usepackage{xwatermark}

%\DeclareTextFontCommand{\texthv}{\fontencoding{T1}\fontfamily{cmss}\fontseries{bx}\fontshape{n}\selectfont}
%\DeclareTextFontCommand{\textcmss}{\usefont{OT1}{cmss}{bx}{n}}
\newcommand{\HRule}{\rule{\linewidth}{0.5mm}}

%\usepackage{tocloft}
%\cftsetindents{section}{0em}{3em}
%\cftsetindents{subsection}{3em}{4em}

%%%%%%%%%%%%%%%%%%%%%%%%%%%%%%%%%%%%%%%%%%%%%%%%%%%%%%%%%%%
%% FUENTES
%%%%%%%%%%%%%%%%%%%%%%%%%%%%%%%%%%%%%%%%%%%%%%%%%%%%%%%%%%%
\usepackage[T1]{fontenc}
% si se usa codificacion T1 y se prefiere
% la fuente textsf en negrita y mas gruesa,
% cargar el paquete lmodern, de lo contrario
% se usa otra fuente
% http://tex.stackexchange.com/questions/227063/fontenc-changes-sans-serif-bold-font-in-koma-script
\usepackage{lmodern}
% si se usa tgpagella esto es solo para cargar las 
% fuentes matematicas
%\usepackage[sc]{mathpazo}
%\linespread{1.05}

%\usepackage{tgpagella}

%\usepackage[minionint,textlf]{MinionPro}
%\renewcommand{\sfdefault}{Myriad-LF}


%\usepackage{newtxtext} % fuente times, texto
%%\usepackage{newtxmath} % fuente times, math
%\usepackage[varg,cmintegrals, cmbraces]{newtxmath}
%\usepackage[cmintegrals, cmbraces, uprightGreek]{newtxmath}

%\usepackage{helvet}
%\renewcommand{\familydefault}{\sfdefault}

\usepackage{microtype}

%%%%%%%%%%%%%%%%%%%%%%%%%%%%%%%%%%%%%%%%%%%%%%%%%%%%%%%%%%
%% LISTADOS DE CODIGO
%%%%%%%%%%%%%%%%%%%%%%%%%%%%%%%%%%%%%%%%%%%%%%%%%%%%%%%%%%%

\usepackage{listings}
\lstset{breaklines=true,basicstyle=\small\ttfamily,} 


%%%%%%%%%%%%%%%%%%%%%%%%%%%%%%%%%%%%%%%%%%%%%%%%%%%%%%%%
% HACKS
%%%%%%%%%%%%%%%%%%%%%%%%%%%%%%%%%%%%%%%%%%%%%%%%%%%%%%%%%%%

% Corrige la separacion entre el numero de la figura/tabla y
% el texto en el indice cuando los numeros son muy
% largos (ej: 10.10), http://tex.stackexchange.com/a/22984
\makeatletter
\renewcommand*\l@figure{\@dottedtocline{1}{1.5em}{2.8em}}% 3em instead of 2.3em
\let\l@table\l@figure
\makeatother

% numerar ecuaciones como chapter.numero
\numberwithin{equation}{chapter}
\numberwithin{figure}{chapter}
\numberwithin{table}{chapter}
%\renewcommand{\figurename}{Figura}
%\renewcommand{\tablename}{Tabla}


%%%%%%%%%%%%%%%%%%%%%%%%%%%%%%%%%%%%%%%%%%%%%%%%%%%%%%%%%%%%
%%              COMANDOS PROPIOS
%%%%%%%%%%%%%%%%%%%%%%%%%%%%%%%%%%%%%%%%%%%%%%%%%%%%%%%%%%%%
\makeatletter
\let\Title\@title
\let\Author\@author
%
\newcommand{\ProfGuia}[1]{\renewcommand\@ProfGuia{#1}}
\newcommand\@ProfGuia{}

\newcommand\@ComisionA{}
\newcommand\@ComisionB{}
\newcommand{\ComisionA}[1]{\renewcommand\@ComisionA{#1\\}}
\newcommand{\ComisionB}[1]{\renewcommand\@ComisionB{#1\\}}
\makeatother

\newcommand{\vect}[1]{\mathbf{#1}}
\newcommand{\vects}[1]{\boldsymbol{#1}}



%%%%%%%%%%%%%%%%%%%%%%%%%%%%%%%%%%%%%%%%%%%%%%%%%%%%%%%%%%%

%\usepackage[pagebackref=true]{hyperref}
\usepackage{hyperref}
\hypersetup{
%    bookmarks=true,         % show bookmarks bar?
    unicode=true,          % non-Latin characters in Acrobat’s bookmarks
%    pdftoolbar=true,        % show Acrobat’s toolbar?
%    pdfmenubar=true,        % show Acrobat’s menu?
    pdffitwindow=false,     % window fit to page when opened
    pdfstartview={FitH},    % fits the width of the page to the window
    pdftitle={Proyecto de Memoria de Título: \Title},    % title
    pdfauthor={\Author},     % author
%    pdfsubject={Subject},   % subject of the document
%    pdfcreator={Creator},   % creator of the document
%    pdfproducer={Producer}, % producer of the document
%    pdfkeywords={keyword1} {key2} {key3}, % list of keywords
    pdfnewwindow=true,      % links in new window
%    colorlinks=false,       % false: boxed links; true: colored links
%    linkcolor=tableau10_10,          % color of internal links (change box color with linkbordercolor)
%    linkbordercolor=tableau10_10,
%    citecolor=green,        % color of links to bibliography
%    filecolor=magenta,      % color of file links
%    urlcolor=cyan           % color of external links
%%%%%%%%%%
    colorlinks=\colorlinks,       % false: boxed links; true: colored links
    linkcolor=tableau10_1,          % color of internal links (change box color with linkbordercolor)
    linkbordercolor=tableau10_10,
    citecolor=tableau10_3,        % color of links to bibliography
    filecolor=magenta,      % color of file links
    urlcolor=tableau10_1           % color of external links
}
